\chapter{Regularization and Structure}

\section{Norm-Based Regularization}

In previous chapters, we saw that highly expressive models can fit training data extremely well, yet risk overfitting. Regularization provides a systematic way to control model complexity and improve generalization.

\medskip

\noindent
\textbf{Basic idea.}  
Augment the empirical risk with a penalty on the parameters:
\[
\min_{\theta \in \Theta} \; \hat{R}_n(\theta) + \lambda \, \Omega(\theta),
\]
where:
\begin{itemize}
    \item $\Omega(\theta)$ measures the complexity of the model,
    \item $\lambda > 0$ controls the strength of regularization.
\end{itemize}

\medskip

\noindent
\textbf{Common choices of penalties.}

\begin{itemize}
    \item \textbf{$\ell_2$ regularization (weight decay):}
    \[
    \Omega(\theta) = \|\theta\|_2^2.
    \]

    \item \textbf{$\ell_1$ regularization:}
    \[
    \Omega(\theta) = \|\theta\|_1.
    \]
\end{itemize}

\medskip

\noindent
\textbf{Interpretation.}
\begin{itemize}
    \item $\ell_2$ regularization favors small, distributed parameters
    \item $\ell_1$ regularization promotes sparsity
\end{itemize}

\medskip

\noindent
\textbf{Geometric viewpoint.}  
Regularization restricts the solution to lie in a region of parameter space with controlled norm.

\medskip

\noindent
\textbf{Statistical viewpoint.}  
Norm penalties can be interpreted as imposing prior distributions on parameters:
\begin{itemize}
    \item $\ell_2$ $\rightarrow$ Gaussian prior
    \item $\ell_1$ $\rightarrow$ Laplace prior
\end{itemize}

\medskip

\noindent
\textbf{Practical effect.}  
Regularization reduces variance and improves stability, often at the cost of a small increase in bias.

\section{Constraints and Penalization}

Regularization can be formulated in two equivalent ways:
\begin{itemize}
    \item \textbf{Penalized form:}
    \[
    \min_\theta \; \hat{R}_n(\theta) + \lambda \Omega(\theta),
    \]
    \item \textbf{Constrained form:}
    \[
    \min_\theta \; \hat{R}_n(\theta) \quad \text{subject to} \quad \Omega(\theta) \leq C.
    \]
\end{itemize}

\medskip

\noindent
\textbf{Insight.}  
These formulations are closely related: for appropriate choices of $\lambda$ and $C$, they yield similar solutions.

\medskip

\noindent
\textbf{Examples.}
\begin{itemize}
    \item Ridge regression: $\ell_2$ penalty
    \item Lasso: $\ell_1$ penalty
\end{itemize}

\medskip

\noindent
\textbf{Interpretation.}  
The constraint limits the effective search space, preventing overly complex solutions.

\medskip

\noindent
\textbf{Dual perspective.}  
Penalization trades off between data fit and model simplicity, while constraints enforce a strict bound on complexity.

\medskip

\noindent
\textbf{Engineering perspective.}  
Penalized formulations are often easier to optimize, while constrained formulations provide clearer geometric intuition.

\section{Implicit Regularization}

Not all regularization is explicit. In many modern learning systems, the optimization algorithm itself induces a preference for certain solutions.

\medskip

\noindent
\textbf{Definition (informal).}  
Implicit regularization refers to the tendency of an algorithm to favor specific solutions even in the absence of explicit penalties.

\medskip

\noindent
\textbf{Example (linear models).}  
In overparameterized linear regression, gradient descent converges to the minimum-norm solution among all interpolating solutions.

\medskip

\noindent
\textbf{Example (deep learning).}
\begin{itemize}
    \item SGD tends to find solutions with particular structure
    \item Different optimizers can lead to different generalization behavior
\end{itemize}

\medskip

\noindent
\textbf{Sources of implicit regularization.}
\begin{itemize}
    \item Initialization
    \item Optimization dynamics
    \item Stochasticity (e.g., mini-batching)
\end{itemize}

\medskip

\noindent
\textbf{Key insight.}  
Even when a model has enough capacity to perfectly fit the data, the training procedure determines \emph{which} solution is found.

\medskip

\noindent
\textbf{Implication.}  
Understanding generalization requires analyzing not only the model and loss, but also the optimization process.

\section{Early Stopping and Practical Methods}

Regularization can also be achieved through practical training strategies.

\medskip

\noindent
\textbf{Early stopping.}  
Instead of minimizing training loss completely, we stop training when validation performance begins to degrade.

\medskip

\noindent
\textbf{Interpretation.}
\begin{itemize}
    \item Early iterations capture large-scale structure
    \item Later iterations may fit noise
\end{itemize}

\medskip

\noindent
\textbf{Connection to regularization.}  
Early stopping behaves similarly to $\ell_2$ regularization in many settings.

\medskip

\noindent
\textbf{Other practical methods.}

\begin{itemize}
    \item \textbf{Dropout:} randomly deactivate units during training
    \item \textbf{Data augmentation:} expand dataset through transformations
    \item \textbf{Noise injection:} add noise to inputs or parameters
\end{itemize}

\medskip

\noindent
\textbf{Interpretation.}
\begin{itemize}
    \item Dropout encourages redundancy and robustness
    \item Data augmentation encodes invariances
    \item Noise smooths the learned function
\end{itemize}

\medskip

\noindent
\textbf{Engineering perspective.}  
These methods are often simple to implement yet highly effective, especially in deep learning systems.

\section{Structure as Regularization}

Beyond explicit penalties and training tricks, model structure itself acts as a powerful form of regularization.

\medskip

\noindent
\textbf{Examples.}
\begin{itemize}
    \item Convolutional layers enforce locality
    \item Residual connections stabilize deep networks
    \item Attention mechanisms capture relational structure
\end{itemize}

\medskip

\noindent
\textbf{Key idea.}  
Structure restricts the class of functions that can be represented, thereby guiding learning toward meaningful solutions.

\medskip

\noindent
\textbf{Insight.}  
In modern machine learning, architecture design is often more important than explicit regularization.

\section{A Unifying Perspective}

Regularization can be understood as controlling the effective complexity of the learned function.

\medskip

\noindent
\textbf{Forms of regularization.}
\begin{itemize}
    \item Explicit (penalties, constraints)
    \item Implicit (optimization dynamics)
    \item Structural (model architecture)
\end{itemize}

\medskip

\noindent
\textbf{Unified objective.}
\[
\text{Fit the data} \quad + \quad \text{control complexity}.
\]

\medskip

\noindent
\textbf{Key components.}
\begin{itemize}
    \item Model class
    \item Regularization mechanism
    \item Optimization algorithm
\end{itemize}

\medskip

\noindent
These elements interact to determine both training performance and generalization.

\section{Outlook}

This chapter introduced regularization as a central tool for controlling model complexity:
\begin{itemize}
    \item norm-based penalties,
    \item constraints and penalization,
    \item implicit effects of optimization,
    \item practical techniques such as early stopping and dropout.
\end{itemize}

\medskip

In the next part of the book, we turn to concrete supervised learning models, beginning with linear regression and least squares.

\medskip

\noindent
\textbf{Guiding principle.}  
Successful learning requires not only fitting data, but shaping the solution space so that meaningful patterns emerge.